% ═══════════════════════════════════════════════════════════════
% LH-03: Elektrische Leistung & Energie - LEHRERHINWEISE
% Physik Klasse 8 | Greenhouse Schools Rostock
% ═══════════════════════════════════════════════════════════════

\documentclass[11pt, a4paper]{article}

\usepackage[utf8]{inputenc}
\usepackage[T1]{fontenc}
\usepackage[ngerman]{babel}

\usepackage{helvet}
\renewcommand{\familydefault}{\sfdefault}

\usepackage{microtype}
\usepackage[margin=1.5cm, top=2cm, bottom=1.5cm]{geometry}
\usepackage{enumitem}
\usepackage{tabularx}
\usepackage{booktabs}
\usepackage{array}
\usepackage{multicol}

\usepackage{amsmath, amssymb}
\usepackage{siunitx}
\sisetup{locale = DE, per-mode = fraction, fraction-function = \dfrac}

\usepackage{fancyhdr}
\usepackage{lastpage}
\usepackage{needspace}

\usepackage{xcolor}
\usepackage{tcolorbox}
\tcbuselibrary{skins}

% ═══════════════════════════════════════════════════════════════
% FARBEN
% ═══════════════════════════════════════════════════════════════

\definecolor{physik}{HTML}{2c5aa0}
\definecolor{hellgrau}{HTML}{F5F5F5}
\definecolor{mittelgrau}{HTML}{888888}
\definecolor{lehrer}{HTML}{7b1fa2}
\definecolor{warnung}{HTML}{b35c00}

\colorlet{fachfarbe}{physik}

% ═══════════════════════════════════════════════════════════════
% VARIABLEN
% ═══════════════════════════════════════════════════════════════

\newcommand{\fach}{Physik}
\newcommand{\thema}{Elektrische Leistung \& Energie}
\newcommand{\klasse}{8}
\newcommand{\stunde}{03}

% ═══════════════════════════════════════════════════════════════
% KOPF- UND FUSSZEILE
% ═══════════════════════════════════════════════════════════════

\pagestyle{fancy}
\fancyhf{}
\renewcommand{\headrulewidth}{0.4pt}
\renewcommand{\footrulewidth}{0pt}
\lhead{\small \fach{} Klasse \klasse{} | \thema{} | \textcolor{lehrer}{\textbf{LEHRERHINWEISE}}}
\rhead{}
\cfoot{\small\textcolor{mittelgrau}{Seite \thepage{} von \pageref{LastPage}}}

% ═══════════════════════════════════════════════════════════════
% BOXEN
% ═══════════════════════════════════════════════════════════════

\newtcolorbox{infobox}[1][]{
    colback=hellgrau,
    colframe=hellgrau,
    boxrule=0pt,
    arc=0pt,
    left=8pt, right=8pt, top=8pt, bottom=6pt,
    borderline north={2pt}{0pt}{fachfarbe},
    before skip=6pt, after skip=6pt,
    #1
}

\newtcolorbox{warnbox}{
    colback=white,
    colframe=warnung,
    boxrule=1pt,
    arc=0pt,
    left=8pt, right=8pt, top=6pt, bottom=6pt,
    before skip=6pt, after skip=6pt
}

\newtcolorbox{lehrerbox}{
    colback=white,
    colframe=lehrer,
    boxrule=1pt,
    arc=0pt,
    left=8pt, right=8pt, top=6pt, bottom=6pt,
    borderline west={3pt}{0pt}{lehrer},
    before skip=6pt, after skip=6pt
}

% ═══════════════════════════════════════════════════════════════
% BEFEHLE
% ═══════════════════════════════════════════════════════════════

\newcommand{\ein}[1]{\,\mathrm{#1}}

\setlength{\parindent}{0pt}
\setlength{\parskip}{6pt}

% ═══════════════════════════════════════════════════════════════
% DOKUMENT
% ═══════════════════════════════════════════════════════════════

\begin{document}

% HEADER
\begin{center}
    {\Large\bfseries\textcolor{lehrer}{Lehrerhinweise: \thema}}\\[0.2em]
    {\small \fach{} Klasse \klasse{} | Stunde \stunde{} | 80 Min (inkl. 20 Min MT)}
\end{center}
\vspace{-0.3em}
\hrule
\vspace{0.8em}

% ═══════════════════════════════════════════════════════════════
% RAHMENPLAN
% ═══════════════════════════════════════════════════════════════

\begin{infobox}
\textbf{Rahmenplan-Bezug (RP Physik MV Gym/Ges 2022, S. 29)}

\textbf{System – Stromkreise:}
\begin{itemize}[leftmargin=1.5em, itemsep=2pt]
    \item Elektrische Leistung: $P_{\text{El}} = U \cdot I$ \hfill [Klasse 7]
    \item Elektrische Energie: $E_{\text{El}} = P_{\text{El}} \cdot t$ \hfill [Klasse 7]
\end{itemize}

\textbf{Hinweis:} Diese Inhalte wurden in Klasse 7 eingeführt und werden hier im Kontext der Stromkreise angewandt und vertieft.
\end{infobox}

% ═══════════════════════════════════════════════════════════════
% STUNDENVERLAUF
% ═══════════════════════════════════════════════════════════════

\section*{Stundenverlauf (80 Min)}

\renewcommand{\arraystretch}{1.3}
\begin{tabularx}{\linewidth}{|c|l|X|l|}
\hline
\textbf{Min} & \textbf{Phase} & \textbf{Aktivität} & \textbf{Material} \\
\hline
0–5 & Einstieg & Stromrechnung zeigen: „Was kostet Duschen?" & Bild \\
\hline
5–12 & Erarbeitung I & Definition $P = \frac{E}{t}$, Berechnung $P = U \cdot I$ & LP Step 1–3 \\
\hline
12–20 & Übung I & ★ Kopfrechnen: $P = 10\ein{V} \cdot 2\ein{A}$ & LP interaktiv \\
\hline
20–28 & Erarbeitung II & Energie $E = P \cdot t$ (aus Definition), Einheiten & LP Step 4–5 \\
\hline
28–35 & Kernfrage & „Warum kWh?" – SI-Einheit Joule vs. haushaltsüblich kWh & LP Step 5–6 \\
\hline
35–45 & \textbf{Praxis} & Labels im Raum suchen, 3 Beispiele notieren & AB Aufg. 1 \\
\hline
45–55 & Anwendung & Stromkosten: Föhn, Durchlauferhitzer & AB Aufg. 2–4 \\
\hline
55–60 & Sicherung & Zusammenfassung Merksätze & AB \\
\hline
60–80 & \textbf{MT} & Stundenleistung (20 Min), 48 BE & MT digital \\
\hline
\end{tabularx}
\renewcommand{\arraystretch}{1.0}

% ═══════════════════════════════════════════════════════════════
% DIFFERENZIERUNG
% ═══════════════════════════════════════════════════════════════

\section*{Differenzierung}

\begin{tabularx}{\linewidth}{|l|X|}
\hline
\textbf{★ Basis} & Direkt einsetzen ($P = U \cdot I$), einfache Zahlen, Label lesen \\
\hline
\textbf{★★ Standard} & Formel umstellen, Einheiten umrechnen (Wh ↔ kWh), Kosten berechnen \\
\hline
\textbf{★★★ Erweitert} & Jahreskosten Standby, Label-Analyse (VA vs. W), Begründung kWh \\
\hline
\end{tabularx}

% ═══════════════════════════════════════════════════════════════
% HÄUFIGE FEHLER
% ═══════════════════════════════════════════════════════════════

\section*{Häufige Fehler \& Reaktionen}

\begin{warnbox}
\textbf{1. Verwechslung Wh und kWh}

\textit{Fehler:} SuS vergessen, durch 1000 zu teilen.

\textit{Reaktion:} „kilo" = 1000 betonen. Eselsbrücke: kWh ist die „große" Einheit für große Energiemengen (Haushalt).
\end{warnbox}

\begin{warnbox}
\textbf{2. Minuten statt Stunden}

\textit{Fehler:} $E = 20\ein{kW} \cdot 15 = 300\ein{kWh}$ (falsch!)

\textit{Reaktion:} Immer Einheiten prüfen! 15 Min = 0,25 h (nicht 15!).
\end{warnbox}

\begin{warnbox}
\textbf{3. ct vs. € verwechselt}

\textit{Fehler:} SuS schreiben „4500 €" statt „4500 ct = 45 €".

\textit{Reaktion:} Strompreis ist in $\frac{\text{ct}}{\text{kWh}}$ → Ergebnis zuerst in ct, dann durch 100 für €.
\end{warnbox}

% ═══════════════════════════════════════════════════════════════
% PRAXISPHASE
% ═══════════════════════════════════════════════════════════════

\section*{Praxisphase: Labels im Raum}

\begin{lehrerbox}
\textbf{Vorbereitung:} Geräte mit sichtbarem Typenschild bereitstellen.

\textbf{Geeignete Geräte im Physikraum:}
\begin{itemize}[leftmargin=1.5em, itemsep=2pt]
    \item Netzgeräte (Input: 230 V, Output: variable V/A)
    \item iPad-/Handy-Ladegeräte (5 V, 2 A oder 9 V, 2 A)
    \item Multimeter (oft 9 V Batterie, mA)
    \item Neonröhren/LED-Panels (W-Angabe)
    \item Overheadprojektor, Beamer (W-Angabe)
\end{itemize}

\textbf{Erwartete Erkenntnisse:}
\begin{itemize}[leftmargin=1.5em, itemsep=2pt]
    \item Nicht alle Geräte haben alle Angaben (V, A, W)
    \item Manche zeigen VA (Scheinleistung) statt W
    \item Frequenz 50 Hz nur bei Wechselstrom
\end{itemize}
\end{lehrerbox}

% ═══════════════════════════════════════════════════════════════
% KOPFRECHENBARE WERTE
% ═══════════════════════════════════════════════════════════════

\section*{Kopfrechenbare Werte (Referenz)}

\begin{multicols}{2}
\textbf{Spannung U:} 10, 12, 20, 230 V

\textbf{Stromstärke I:} 0,5, 1, 2, 3, 5, 10 A

\textbf{Leistung P:} 20, 50, 100, 500, 1000, 2000 W

\textbf{Zeit t:} 10 Min, 30 Min, 1 h, 2 h, 5 h

\textbf{Strompreis:} $\mathbf{30\,\frac{ct}{kWh}}$ (gerundet, realistisch)

\columnbreak

\textbf{Typische Rechnungen:}

$P = 12\ein{V} \cdot 4\ein{A} = 48\ein{W}$

$E = 2000\ein{W} \cdot 0,5\ein{h} = 1000\ein{Wh} = 1\ein{kWh}$

Kosten: $1\ein{kWh} \cdot 30\ein{ct} = 30\ein{ct}$
\end{multicols}

% ═══════════════════════════════════════════════════════════════
% MATERIAL
% ═══════════════════════════════════════════════════════════════

\section*{Material-Checkliste}

\begin{itemize}[leftmargin=1.5em, itemsep=3pt]
    \item[$\square$] LP-03-elektrische-leistung.html (Beamer/Tablets)
    \item[$\square$] AB-03-elektrische-leistung.pdf (Klassensatz)
    \item[$\square$] MT-03-elektrische-leistung.html (Tablets/PCs)
    \item[$\square$] Geräte mit Typenschild (mind. 5 verschiedene)
    \item[$\square$] Klemmbretter für Label-Suche
    \item[$\square$] Optional: Echte Stromrechnung (anonymisiert)
\end{itemize}

% ═══════════════════════════════════════════════════════════════
% MT-HINWEISE
% ═══════════════════════════════════════════════════════════════

\section*{MT-03: Hinweise zur Auswertung}

\begin{infobox}
\textbf{Aufbau:} 48 BE, 20 Minuten, 3 Niveaus

\begin{tabular}{llr}
\textbf{Teil I} & Grundlagen (★) & 25 BE \\
\textbf{Teil II} & Anwendung (★★) & 13 BE \\
\textbf{Teil III} & Transfer (★★★) & 10 BE \\
\end{tabular}

\textbf{Differenzierte Wertung:}
\begin{itemize}[leftmargin=1.5em, itemsep=2pt]
    \item ★ Basis: Nur Teil I bearbeiten (max. 25 BE)
    \item ★★ Standard: Teil I + II bearbeiten (max. 38 BE)
    \item ★★★ Erweitert: Alle Teile (max. 48 BE)
\end{itemize}

\textbf{III.2 (Freitext):} Manuell bewerten – 5 BE für vollständige Begründung (Größe der Zahl, Handlichkeit, praktische Nutzung).
\end{infobox}

% ═══════════════════════════════════════════════════════════════
% SPRACHLICHE HINWEISE
% ═══════════════════════════════════════════════════════════════

\section*{Sprachliche Hinweise}

\begin{lehrerbox}
\textbf{„Energieverbrauch" → „Energiewandlung"}

Physikalisch wird Energie nicht „verbraucht", sondern umgewandelt. Im Haushalt wird elektrische Energie in Wärme, Licht etc. umgewandelt.

\textit{Didaktischer Kompromiss:} Der Begriff „Verbrauch" ist alltagssprachlich etabliert. Kurz thematisieren, aber nicht überbetonen.
\end{lehrerbox}

\end{document}
